\section{RQ3: Models, Engines, Benchmarks, Faithfulness and Reproducibility}
\label{sec:evaluation}

\begin{table}[t]
\centering
\caption{Language-model families, symbolic engines and benchmarks named in the abstracts of included studies (open-vocabulary extraction by two coders; a study counts if either coder listed it).}
\label{tab:entities}
\footnotesize
\begin{tabular}{@{}lr@{}}
\toprule
Named in abstract & Studies \\
\midrule
\multicolumn{2}{@{}l}{\textit{Language-model families (199 studies name one)}} \\
\quad GPT-4 family & 93 \\
\quad GPT-3/3.5, ChatGPT, Codex & 47 \\
\quad Llama family & 46 \\
\quad DeepSeek family & 29 \\
\quad Gemini/PaLM/Bard & 25 \\
\quad Specialised provers & 15 \\
\quad Claude & 13 \\
\quad BERT family & 11 \\
\quad Qwen family & 11 \\
\midrule
\multicolumn{2}{@{}l}{\textit{Symbolic engines (215 studies name one)}} \\
\quad Python interpreter/executor & 70 \\
\quad Lean & 62 \\
\quad Planners (PDDL) & 19 \\
\quad SAT/SMT (Z3, cvc5) & 17 \\
\quad Optimisation solvers (LP/MILP/CP) & 15 \\
\quad Isabelle & 13 \\
\quad Logic programming (Prolog, ASP, Datalog) & 9 \\
\quad Theorem provers (FOL) & 8 \\
\quad Dafny/Verus/Why3/F* & 8 \\
\midrule
\multicolumn{2}{@{}l}{\textit{Benchmarks (353 studies name one)}} \\
\quad miniF2F & 42 \\
\quad MATH & 16 \\
\quad ProofNet & 14 \\
\quad GSM8K & 12 \\
\quad ALFWorld & 8 \\
\quad PutnamBench & 8 \\
\quad HumanEval & 7 \\
\quad AIME & 6 \\
\quad WikiTableQuestions & 5 \\
\bottomrule
\end{tabular}

\end{table}

\textbf{What abstracts name.} Table~\ref{tab:entities} summarises the entities named in abstracts. The two coders agreed well on model families and benchmarks (mean Jaccard \cnt{jaclm} and \cnt{jacbenchmark}) and less well on engines (\cnt{jacengine}), so engine counts are approximate. Only \cnt{lmstudies} studies (\cnt{lmstudiespct}\%) name a model family in the abstract; the counts describe reporting in abstracts, not usage. Among them, GPT-4-family models are the most named (\cnt{lmGPT4}). Proprietary and open-weight families are used side by side: \cnt{lmprop} studies name a proprietary family and \cnt{lmopen} an open-weight one, while \cnt{lmproponly} name only proprietary families. The Python interpreter (\cnt{engPython}) and Lean (\cnt{engLean}) are the most named engines, which reflects the weight of program-aided reasoning and formal theorem proving.

\textbf{Benchmarks are fragmented.} \cnt{benchstudies} studies name at least one benchmark in the abstract, and together they name \cnt{benchdistinct} distinct benchmarks. Of these, \cnt{benchsingle} (\cnt{benchsinglepct}\%) are named by a single study, and only \cnt{benchfiveplus} are named by five or more. miniF2F is the only benchmark named by more than 20 studies (\cnt{benchminif2f}). This conclusion no longer depends on a fixed dictionary, because the vocabulary is open. Abstracts under-report benchmarks, though, so the true overlap is somewhat higher.

% miniF2F pass rates reported in abstracts of included studies (data/derived/plot_minif2f.csv).
\begin{figure}[t]
\centering
\begin{tikzpicture}
\begin{axis}[
  width=\columnwidth, height=52mm,
  xmin=2021.6, xmax=2026.5, ymin=20, ymax=105,
  xtick={2022,2023,2024,2025,2026},
  xticklabel style={font=\scriptsize, /pgf/number format/1000 sep={}},
  yticklabel style={font=\scriptsize},
  ylabel={Reported miniF2F pass rate (\%)}, ylabel style={font=\scriptsize},
  axis lines*=left, ymajorgrids, grid style={black!10},
  legend style={font=\tiny, at={(0.02,0.98)}, anchor=north west, draw=none, fill=white, fill opacity=0.85, text opacity=1},
  legend cell align=left,
  scatter/classes={Lean={mark=*, nsblue}, Isabelle={mark=triangle*, A4c}, unspecified={mark=square*, nsgrey}},
]
\addplot[scatter, only marks, scatter src=explicit symbolic, mark size=2pt]
  table[col sep=comma, x expr={\thisrow{year}+0.09*(mod(\coordindex,5)-2)}, y=score, meta=lang] {data/derived/plot_minif2f.csv};
\legend{Lean, Isabelle, not stated}
\end{axis}
\end{tikzpicture}
\caption{Headline miniF2F results stated in the abstracts of included studies (\cnt{mfcomparable} absolute pass rates; \cnt{mfreported} abstracts give a miniF2F number, the rest are gains or variant benchmarks). The best reported rate rises from \cnt{mfmax2022}\% (2022) to \cnt{mfmax2026}\% (2026). Points are not directly comparable: sampling budgets range from pass@1 to thousands of attempts, and splits and formal languages differ. Horizontal jitter separates points within a year.}
\label{fig:minif2f}
\end{figure}


\textbf{A results-level view of the one shared benchmark.} Of \cnt{mfmention} D1 or miniF2F-mentioning studies, \cnt{mfreported} state a miniF2F number in the abstract and \cnt{mfcomparable} state an absolute pass rate (Fig.~\ref{fig:minif2f}). The best reported rate rises from \cnt{mfmax2022}\% in 2022~\cite{wu2022autoformalization} to \cnt{mfmax2025}\% in 2025~\cite{ospanov2025apollo} and \cnt{mfmax2026}\% in 2026~\cite{varambally2026hilbert}. These numbers are not a leaderboard. Sampling budgets range from a single attempt to thousands, splits and formal languages differ, and audits find mis-specified statements and evaluation loopholes in the benchmark itself~\cite{ospanov2025minif2f,ammanamanchi2026faults}. The figure shows two things: the most shared benchmark of the field is close to saturation, and results are reported in a way that prevents direct comparison.

\textbf{Formalisation faithfulness is rarely reported.} For A1 and A2 studies whose model writes formal statements, correctness depends on whether the formal statement means what the problem says. Of \cnt{faithn} such studies, \cnt{faithyes} (\cnt{faithyespct}\%) report the accuracy or faithfulness of the formalisation separately from end-task accuracy in the abstract, \cnt{faithno} do not, and \cnt{faithunc} are unclear. Reporting is most common in D1 (\cnt{faithD1yes} of \cnt{faithD1}), where autoformalisation is itself a task~\cite{li2024autoformalize,liu2025rethinking,chen2026reform}. It is almost absent in logical reasoning (\cnt{faithD2yes} of \cnt{faithD2}) and mathematical problem solving (\cnt{faithD7yes} of \cnt{faithD7}), where translate-and-solve pipelines are common. This was coded from abstracts, so some studies report faithfulness only in the body; the figure is a lower bound on reporting, not on practice.

\begin{table}[t]
\centering
\caption{Release of author code, models or data, verified on the web: a simple random sample of included studies by architecture, and a census of A7 studies.}
\label{tab:artifacts}
\footnotesize
\setlength{\tabcolsep}{4pt}
\begin{tabular}{@{}lrrrr@{}}
\toprule
Stratum & $n$ & Released & Not found & Unclear \\
\midrule
\quad A1 Translate-and-solve & 23 & 17 & 5 & 1 \\
\quad A2 Checker-in-the-loop & 59 & 39 & 12 & 8 \\
\quad A3 Symbolic guidance & 14 & 8 & 5 & 1 \\
\quad A4 LM-induced artifacts & 10 & 7 & 3 & 0 \\
\quad A5 Symbolic-to-neural transfer & 14 & 12 & 2 & 0 \\
\quad A6 Symbolic-oracle evaluation & 25 & 23 & 2 & 0 \\
\quad A7 LM as interface & 5 & 5 & 0 & 0 \\
\midrule
Random sample & 150 & 111 & 29 & 10 \\
A7 census & 22 & 14 & 3 & 5 \\
\bottomrule
\end{tabular}

\end{table}

\textbf{Artifacts are usually released.} In the random sample of \cnt{artrandn} studies, \cnt{artrandyes} have an author-released code, model or data artifact, \cnt{artrandno} have none that could be found and \cnt{artrandunc} could not be determined (Table~\ref{tab:artifacts}). The estimated release rate is \cnt{artrandrate}\% (95\% CI \cnt{artrandlo}--\cnt{artrandhi}). It is \cnt{artrandknown}\% among determinable studies, and \cnt{artrandupper}\% if all unclear studies released. In the A7 census, \cnt{artasevenyes} of \cnt{artasevenn} studies release an artifact. Some groups of sampled studies are small, so per-architecture differences should not be over-read. The pilot estimated release from links in PDF text and reported a markedly lower rate, which shows how much that method under-counts.
